% Claude 1962 proper study, Seventeenth Sunday after Pentecost: leaf-local
% presentation. This file owns typography only; every textual, historical,
% exegetical and rights judgment lives in the section files and in the
% leaf's research records. It is shared by the three entrypoints of the leaf
% (main.tex, synthesis.tex, homily.tex) so that a component used by more than
% one of them sets identically in each. Because it is shared, it is an input to
% each of the three reviews: an environment added here for one entrypoint
% re-opens the other two documents' accepted reviews, and the study owns the
% file whichever entrypoint the addition serves.
%
% Web-edition constraint: the canonical study is declared web-eligible, so no
% landscape, multicols, TikZ or rule-line construct is used here or in any
% research component.

% --- Type and measure -------------------------------------------------------
% Palatino (TeX Gyre Pagella through newpxtext) at 11pt on 15pt leading, on a
% measure of about seventy characters.
\usepackage{newpxtext}
\linespread{1.1}
\geometry{letterpaper,left=1.3in,right=1.3in,top=1in,bottom=1.05in,
  headheight=14pt,headsep=20pt,footskip=32pt}
\setlength{\parindent}{1.2em}
\setlength{\parskip}{0pt plus 1pt}
\setlength{\emergencystretch}{2em}
\widowpenalty=10000
\clubpenalty=10000
\displaywidowpenalty=10000
\brokenpenalty=10000
\usepackage{xurl}
\newcommand{\work}[1]{\textit{#1}}

% Unnumbered sectioning that still reaches the contents and the running head.
% The bookmark package gives each unnumbered heading its own outline anchor.
\usepackage{bookmark}
\setcounter{secnumdepth}{0}
\setcounter{tocdepth}{1}
\usepackage{titlesec}
\titleformat{\section}{\normalfont\Large\bfseries\raggedright}{}{0pt}{}
\titlespacing*{\section}{0pt}{3.2ex plus 1ex minus .2ex}{1.4ex plus .2ex}
\titleformat{\subsection}{\normalfont\large\bfseries\raggedright}{}{0pt}{}
\titlespacing*{\subsection}{0pt}{2.6ex plus .8ex minus .2ex}{0.9ex plus .2ex}
\titleformat{\subsubsection}{\normalfont\normalsize\bfseries\itshape\raggedright}{}{0pt}{}
\titlespacing*{\subsubsection}{0pt}{2ex plus .6ex minus .2ex}{0.6ex plus .1ex}

% Quiet running matter: the Sunday on the left, the current section on the
% right, the folio centred below.
%
% Width discipline. fancyhdr laps each of the two head strings out of a parbox
% the full width of the head, so the left and the right string are free to
% occupy the same space: a right string longer than what the left one leaves
% simply prints over it, the \hss glue absorbs the overflow, and the log
% reports neither an Overfull box nor anything else. That is a silent defect
% and it happened here. The right box is therefore cut to what remains of the
% head after the fixed left string and a gutter, so that a string too long for
% it wraps inside its own box rather than reaching the left string, and the
% second line trips fancyhdr's \headheight warning, which a log inspection
% catches. The gutter is real white space and not glue: no setting of the two
% strings can close it.
%
% The discipline bounds the damage; it is not the whole remedy. A wrapped head
% is legible but ugly, so every section whose title is wider than
% \tptheadrightwidth declares a short running form with \runninghead just
% before its heading, and the full title stands unshortened in the section and
% in the contents. \section* headings, which set no mark at all and would
% otherwise leave the previous section's title standing over pages they do not
% own, call \runninghead after the heading for the same reason.
%
% \runninghead sets the mark at once and also hands the same string to the next
% \sectionmark, which prints it in place of the title it is given. Both halves
% are needed. \rightmark takes the FIRST mark on the page, not the last, so a
% mark set only after \section is beaten by the heading's own mark on the one
% page where it matters most --- the page the heading opens; and \section* sets
% no mark at all, so for a starred heading the immediate one is the only one.
% An earlier attempt wrote \markright after the heading and left pages 20 and
% 26 overprinting exactly as they were.
\usepackage{fancyhdr}
\pagestyle{fancy}
\fancyhf{}
\renewcommand{\headrulewidth}{0pt}
% What may stand in \runninghead's body is decided by pandoc, not by LaTeX. The
% web converter does not scan this preamble file, but it does carry its
% \newcommand definitions into the document it hands pandoc, and it expands
% every \runninghead of the body there. \gdef and \markright are safe --- the
% one defines, the other pandoc knows and drops, which is right, because a
% running head is not reader-facing text in a web edition. A \newif boolean is
% not: \global\tptshortheadtrue here failed the conversion outright. The test
% is therefore an \ifx, and it lives in \sectionmark, which LaTeX alone calls
% and pandoc never expands. \let, not an empty \newcommand, sets and restores
% the unset state, so \ifx compares two tokens identical by construction.
\newcommand{\tptnoshorthead}{}
\let\tptshortheadtext\tptnoshorthead
\newcommand{\runninghead}[1]{\gdef\tptshortheadtext{#1}\markright{#1}}
\renewcommand{\sectionmark}[1]{%
  \ifx\tptshortheadtext\tptnoshorthead
    \markright{#1}%
  \else
    \expandafter\markright\expandafter{\tptshortheadtext}%
    \global\let\tptshortheadtext\tptnoshorthead
  \fi}
\newcommand{\tptheadsunday}{Seventeenth Sunday after Pentecost}
\newlength{\tptheadgutter}
\setlength{\tptheadgutter}{1.5em}
\newlength{\tptheadleftwidth}
\newlength{\tptheadrightwidth}
\newcommand{\tptheadright}[1]{%
  \settowidth{\tptheadleftwidth}{\footnotesize\scshape\tptheadsunday}%
  \setlength{\tptheadrightwidth}{%
    \dimexpr\headwidth-\tptheadleftwidth-\tptheadgutter\relax}%
  \parbox[b]{\tptheadrightwidth}{\raggedleft\footnotesize\itshape #1}}
\fancyhead[L]{\footnotesize\scshape\tptheadsunday}
\fancyhead[R]{\tptheadright{\nouppercase{\rightmark}}}
\fancyfoot[C]{\footnotesize\thepage}
\fancypagestyle{plain}{\fancyhf{}\renewcommand{\headrulewidth}{0pt}%
  \fancyfoot[C]{\footnotesize\thepage}}

% Physical-page evidence for the concise companion's fixed opening, whose
% entrypoint shares this file; the study itself sets no page markers.
\usepackage{zref-user,zref-abspage}

% --- Characters the Missal prints outside T1 ---------------------------------
% The versicle sign, as a slashed V, so that the appointed text is set as
% printed; the acute over the ae-ligature (quǽsumus, sǽcula); and the cross
% before the Gospel heading.
\usepackage{graphicx}
\newcommand{\Vsign}{\leavevmode\hbox{\rlap{\kern0.16em\raisebox{0.12ex}{\scalebox{0.85}[1.0]{/}}}V}}
\DeclareUnicodeCharacter{2123}{\Vsign}
\DeclareUnicodeCharacter{01FD}{\'{\ae}}
\DeclareUnicodeCharacter{01FC}{\'{\AE}}

% --- Appointed texts -----------------------------------------------------------
% The Latin of the 1962 editio typica, set as a block; and the English of an
% identified public-domain witness, whose name and locus stand above it.
\newenvironment{latinproper}{\begin{quote}\itshape}{\end{quote}}
\newenvironment{englishwitness}[1]{\begin{quote}\textsc{#1}\par\nopagebreak}{\end{quote}}
\newcommand{\properlocus}[1]{\par{\small\noindent #1\par}}

% --- The four senses of each reading ---------------------------------------------
\newenvironment{fourSenses}{\begin{description}}{\end{description}}

% --- Scriptural Date and Location -----------------------------------------------
% Two-tier dossier table after the 1962 profile's page-2 standard: a four-field
% summary row, then a full-width explanatory row, with an italic event row for
% a narrated event and \midrule only between complete dossiers.
%
% Measurement note. The profile fixes the summary tier at 0.14 / 0.18 / 0.35 /
% 0.16\linewidth and the explanatory row at 0.92\linewidth, and permits a
% change only when the complete inventory cannot remain legible after the
% hierarchy and evidence have been preserved. Two measurements are varied, and
% nothing else the profile fixes:
%
%   D1. Summary-tier widths 0.12 / 0.14 / 0.29 / 0.39\linewidth (0.94 in
%       all, so that the summary tier spans the same measure as the row
%       beneath it). The generated annotation prints six disputed composition
%       labels at the Gospel beside its event statement and two at the
%       Epistle, each in the source's own words; at the profile's
%       0.16\linewidth the Gospel's Date cell alone ran to thirteen lines, and
%       at 0.32\linewidth the sheet still ran onto a second page.
%   D2. Explanatory and event rows at 0.98\linewidth in \scriptsize, so that
%       seven dossiers stand on the single physical page the concise
%       companion's page 2 must occupy. The study's own terminal appendix,
%       which the profile does not confine to one page, prints the same seven
%       dossiers across two.
%
% Kept as the profile fixes them: the two-tier hierarchy, the four summary
% fields, \footnotesize for the summary tier, \arraystretch 1.02, \LTpre
% 0.15em, \LTpost 0pt, 0.1em after ordinary dossier rows, \cmidrule(lr){1-4}
% within a dossier and \midrule between dossiers.
\newenvironment{dossiertable}
{\footnotesize
 \renewcommand{\arraystretch}{1.02}%
 \setlength{\LTpre}{0.15em}\setlength{\LTpost}{0pt}%
 \begin{longtable}{@{}>{\raggedright\arraybackslash}p{0.12\linewidth}>{\raggedright\arraybackslash}p{0.14\linewidth}>{\raggedright\arraybackslash}p{0.29\linewidth}>{\raggedright\arraybackslash}p{0.39\linewidth}@{}}
 \toprule
 \textbf{Proper} & \textbf{Citation} & \textbf{Location} & \textbf{Date}\\
 \midrule
 \endhead}
{\bottomrule\end{longtable}}

% \chronodate{element-key}{content}: one dossier row's Date cell. It typesets
% the content and nothing else. Every biblical date printed in this leaf comes
% from the Scripture chronology corpus through the generated projection in
% research/chronology-annotations.tex, which owns \chronologyannotation.
\newcommand{\chronodate}[2]{#2}

\newcommand{\dossierrowsize}{\scriptsize}
\newcommand{\dossierprose}[1]{\multicolumn{4}{@{}>{\raggedright\arraybackslash\dossierrowsize}p{0.98\linewidth}@{}}{#1}\\[0.1em]}
\newcommand{\dossierevent}[1]{\multicolumn{4}{@{}>{\raggedright\arraybackslash\dossierrowsize}p{0.98\linewidth}@{}}{\textit{#1}}\\}

% --- Comparison tables ----------------------------------------------------------
\newenvironment{comparisontable}[4]
{\small\renewcommand{\arraystretch}{1.18}%
 \begin{longtable}{@{}>{\raggedright\arraybackslash}p{0.17\linewidth}>{\raggedright\arraybackslash}p{0.26\linewidth}>{\raggedright\arraybackslash}p{0.26\linewidth}>{\raggedright\arraybackslash}p{0.26\linewidth}@{}}
 \toprule
 \textbf{#1} & \textbf{#2} & \textbf{#3} & \textbf{#4}\\
 \midrule
 \endhead}
{\bottomrule\end{longtable}}

% The concise companion's first physical page: the map of the appointed
% elements and, beneath it, the four overview rows. Both are set one step
% smaller than the study's map so that the complete inventory and the four
% senses stand together on the single page the concise profile fixes for them,
% and their column widths are chosen so that the two tables rule to exactly the
% same measure: 0.96\linewidth of columns plus two inter-column gaps in the
% three-column map, and 0.9788\linewidth plus one gap in the two-column
% overview, each with \tabcolsep at 4pt, which come to the same width.
\newenvironment{concisemaptable}
{\footnotesize\renewcommand{\arraystretch}{1.04}\setlength{\tabcolsep}{4pt}%
 \begin{longtable}{@{}>{\raggedright\arraybackslash}p{0.15\linewidth}>{\raggedright\arraybackslash}p{0.25\linewidth}>{\raggedright\arraybackslash}p{0.56\linewidth}@{}}
 \toprule
 \textbf{Element} & \textbf{Text and reference} & \textbf{What it says}\\
 \midrule
 \endhead}
{\bottomrule\end{longtable}}

\newenvironment{conciseoverview}
{\footnotesize\renewcommand{\arraystretch}{1.06}\setlength{\tabcolsep}{4pt}%
 \begin{longtable}{@{}>{\bfseries\raggedright\arraybackslash}p{0.15\linewidth}>{\raggedright\arraybackslash}p{0.8288\linewidth}@{}}
 \toprule
 \textbf{Sense} & \textbf{Synthesis}\\
 \midrule
 \endhead}
{\bottomrule\end{longtable}}

\newenvironment{maptable}
{\small\renewcommand{\arraystretch}{1.15}%
 \begin{longtable}{@{}>{\raggedright\arraybackslash}p{0.17\linewidth}>{\raggedright\arraybackslash}p{0.27\linewidth}>{\raggedright\arraybackslash}p{0.52\linewidth}@{}}
 \toprule
 \textbf{Element} & \textbf{Text and reference} & \textbf{What it says}\\
 \midrule
 \endhead}
{\bottomrule\end{longtable}}
